%% ma: L12-B01-0008
%% lop: 12
%% chuong: 1
%% bai: 1
%% dang: 12.1.1
%% loai: TN4
%% muc: TH
%% dap_an: "A"
%% nguon: "(NBV)-ĐỀ SỐ 6-MỖI NGÀY 1 ĐỀ THI 2025.pdf (D06, MỖI NGÀY 1 ĐỀ - Đề số 6), Phần I Câu 1"
%% kien_thuc: "Bài 1 (điều kiện cực trị của hàm số bậc ba: y'(x0)=0 và đổi dấu)"
%% trang_thai: chinh-thuc
%% ghi_chu: "Trùng ý tưởng với dạng 12.1.1 (xác định hàm bậc ba từ cực trị) nên nhập cùng dạng. Đề gốc ghi 'điểm cực trị của hàm số', đã sửa thành 'của đồ thị hàm số' cho đúng thuật ngữ. Đáp án gốc A đúng"
\begin{ex}
Cho biết $A(-1;14)$ là một điểm cực trị của đồ thị hàm số $y=2x^3-9x^2+mx+n$. Khi đó $m+n$ bằng
\choice{\True $-23$}{$23$}{$-11$}{$7$}
\loigiai{$y'=6x^2-18x+m$. Vì $A(-1;14)$ là điểm cực trị nên $y'(-1)=0$ và $y(-1)=14$, tức là $24+m=0$ và $-11-m+n=14$. Suy ra $m=-24$, $n=1$.

Thử lại: $y'=6(x+1)(x-4)$ đổi dấu khi qua $x=-1$ nên $A$ là điểm cực trị. Vậy $m+n=-23$.}
\end{ex}
